% ANCHOR: after_label def:c7-component
\begin{remark}[Physics notation: components as quantum transformations]
With the constant ghosts of $\mathfrak{psl}(4|4)$ added to the BRST differential, the component along $x_1$ is the quantum transformation of the renormalised operator $O$ under the generator $x_1$, the one-loop form of $[x_1,O]$; at order $\hbar^0$ it is the classical variation $\{S^{(1)}(x_1),O\}$. For the odd generators its coefficients along $\xi_\alpha{}^i$, $\bar\xi_{\dot\alpha i}$, $\eta^\alpha{}_i$, $\bar\eta^{\dot\alpha i}$ are the transformations under $Q^i_\alpha$, $\bar Q_{\dot\alpha i}$, $S^\alpha{}_i$, $\bar S^{\dot\alpha i}$ (Definition \ref{def:c8-components}). The parameter $\delta$ is the source of $O$, added to the action as $\delta\,O(\varepsilon)$. A local representative $\int\omega\,f_{x_1}=\int d^4x\,\omega(x)f_{x_1}(x)$ is the smeared composite operator $f_{x_1}(x)$, a polynomial in $A_m$, $A_{ij}$, $\lambda_\alpha{}^i$, $\bar\lambda_{\dot\alpha i}$, $c$, the antifields and their derivatives at $x$.
\end{remark}

% ANCHOR: end_of_section 7.1
\begin{remark}[Physics notation: BRST-exactness of the components]
In the language of the BRST operator extended by constant ghosts (the extended Slavnov--Taylor operator of algebraic renormalisation), Lemma \ref{lem:c7-image} states that the transformation of $O$ under $x_1$ is the $Q$-variation of the part of $O[\infty]$ linear in the constant ghost of $x_1$. Its physical content is carried by the local composite operator $f_{x_1}(x)$, to which the selection rules of this chapter apply.
\end{remark}

% ANCHOR: after_label def:c7-vdot
\begin{remark}[Physics notation: conformal transformations of the fields]
On the fields, $v.$ is the two-component transformation of $S^{(1)}(v)$ written out after Definition \ref{def:c2-S1}: the conformal transformation of the vector $A_m$, of the scalars $A_{ij}$ of dimension $1$, of the Weyl spinors $\lambda_\alpha{}^i$, $\bar\lambda_{\dot\alpha i}$ of dimension $\tfrac32$ and of the ghost $c$ of dimension $0$, with the spin term $-\tfrac18X_{mn}(\sigma^m\bar\sigma^n)_\alpha{}^\beta$ on $\lambda_\alpha{}^i$ and the $SU(4)_R$ rotation $r^i{}_j=\tfrac14R_{IJ}(\bar\Sigma^J\Sigma^I)^i{}_j$, traceless and anti-Hermitian for real $R$. Here $X_{mn}=\partial_mv_n-\partial_nv_m$ is the local rotation and $\tfrac14\operatorname{div}v$ the local scale factor, $\partial_mv_n+\partial_nv_m=\tfrac12\operatorname{div}v\,\delta_{mn}$ (the conformal Killing equation), and for $A_m$ the weight term sits in $A_n\partial_mv^n$. The condition on the antifields makes $A^*$, $\lambda^*$, $c^*$ primaries of dimension $4-w$, i.e. $3$, $\tfrac52$, $4$ (Lemma \ref{lem:c7-weights}), the dimensions for which the field--antifield pairing, of total dimension $4$, is invariant.
\end{remark}

% ANCHOR: after_label prop:c7-cov
\begin{remark}[Physics notation: the covariant generators]
The covariant generators are the translations $P_b$, the rotations $\mathfrak{so}(4)=\mathfrak{su}(2)_L\oplus\mathfrak{su}(2)_R$ (the generators $M$, $\bar M$ of [SW, (8)]) and the R-symmetry $\mathfrak{su}(4)_R\cong\mathfrak{so}(6)_R$ (the generators $R$). They commute with the Feynman-gauge kinetic operator $H=-\partial^2$, with the gauge-fixing operator $Q^{GF}$ and with the proper-time cutoff $e^{-\varepsilon|p|^2}$, so the regulated propagators and the counterterms are Poincaré and $SU(4)_R$ covariant. The spin representation $\tfrac14\rho\circ q$ of a rotation acts on $\lambda_\alpha{}^i$ by $-\tfrac18X_{mn}(\sigma^m\bar\sigma^n)_\alpha{}^\beta$, on $\bar\lambda_{\dot\alpha i}$ from the right by $+\tfrac18X_{mn}(\bar\sigma^m\sigma^n)^{\dot\beta}{}_{\dot\alpha}$, and an R-rotation acts on the index of $W$ by $r^i{}_j$.
\end{remark}

% ANCHOR: after_label cor:c7-schur
\begin{remark}[Physics notation: $\mathbf{20'}$ and Konishi mixing]
Under $SO(6)_R$, $\mathrm{Sym}^2\mathbf6=\mathbf1\oplus\mathbf{20'}$. With $\phi_I=A^a_IT^a$ and $(\iota_{w_1}A,\iota_{w_2}A)=4\hbar\,w_{1I}w_{2J}\operatorname{tr}(\phi_I\phi_J)$, the $\mathbf{20'}$ operators $\operatorname{tr}(\phi_I\phi_J)-\tfrac16\delta_{IJ}\operatorname{tr}(\phi_K\phi_K)$ renormalise multiplicatively, and the Konishi operator $\operatorname{tr}(\phi_I\phi_I)$ is renormalised into the span of itself, of the gauge-variant $(A|_V,A|_V)=2\hbar A^a_mA^a_m$ and of the identity operator $\int\omega$. The coefficient of the identity, $\frac{2\hbar\dim\mathfrak g\,(w_1,w_2)}{16\pi^2\varepsilon}$ with $\dim\mathfrak{su}(N)=N^2-1$, is a power divergence, the subtraction of the vacuum expectation value of the bilinear; the coefficient of $(A|_V,A|_V)$ is $0$.
\end{remark}

% ANCHOR: end_of_section 7.2
\begin{remark}[Physics notation: classical action of $P$, $M$, $\bar M$, $R$]
Remark \ref{rem:c7-internal} in physics language: for BRST-invariant $O$ the one-loop transformations under $P_b$, $M$, $\bar M$ and $R$ can be chosen to be the classical ones (derivative, rotation of the Lorentz indices, $SU(4)_R$ rotation of the $W$ indices), with no anomalous term. The internal insertions exist because the gauge-fixing operator squares to a nonzero operator; in components
\[
(Q^{GF})^2A^*_m=-2\,\partial_mA^*_m\quad(\text{in the slot of }c),
\]
the composition of $A^*_m\mapsto-2A^*_m$ (slot of $A_m$) with $A_m\mapsto\partial_mA_m$ (slot of $c$), equal to $2d^*(A^*|_V)$ with $d^*=-\partial_m$ on one-forms.
\end{remark}

% ANCHOR: after_label prop:c7-dil
\begin{remark}[Physics notation: dilatation and special conformal fields]
In components the dilatation is $v^m=2x^m$, $\operatorname{div}v=8$, and the special conformal field along $b$ is
\[
v^m=x^2b^m-2(b\cdot x)x^m,\qquad X_{mn}=4(x_mb_n-x_nb_m),\qquad \operatorname{div}v=-8\,b\cdot x .
\]
On the scalars $v.A_{ij}=\big(x^2b^n-2(b\cdot x)x^n\big)\partial_nA_{ij}-2(b\cdot x)A_{ij}$, the special conformal transformation $b^n\big(x^2\partial_n-2x_n\,x\cdot\partial-2w\,x_n\big)$ of a scalar primary of dimension $w=1$. The proposition states that $\Delta$ and $K_b$ fail to commute with the Feynman-gauge operator $Q^{GF}$, and that the dilatation acts on the propagator as the rescaling of the proper-time cutoff: $(v.\otimes1+1\otimes v.)P(\varepsilon,\infty)=-4\varepsilon\frac{d}{d\varepsilon}P(\varepsilon,\infty)$ is, on each component, $4\varepsilon\,e^{-\varepsilon|p|^2}$ times the numerator of the propagator.
\end{remark}

% ANCHOR: after_label prop:c7-susy
\begin{remark}[Physics notation: free-field supersymmetry and the gauge-fixing operator]
$Q_1\in S_+\otimes W$ is a combination of the supercharges $Q^i_\alpha$, in $(\mathbf2,\mathbf1,\mathbf4)$, and $Q_1\in S_-\otimes W^*$ of the $\bar Q_{\dot\alpha i}$, in $(\mathbf1,\mathbf2,\bar{\mathbf4})$. The field transformations of $S^{(1)}(Q_1)$ are the $\mathcal N=4$ supersymmetry transformations written out at the end of Section 2.3, $Q_1.A=-\Gamma(Q_1,\lambda)$, $Q_1.\lambda=-\tfrac12\rho(F_A)Q_1$; the operator $Q_1.$ of the proposition is their free part ($F_{mn}\to\partial_mA_n-\partial_nA_m$, $D_mA^{ij}\to\partial_mA^{ij}$, the term $\tfrac14[A^{ik},A_{kj}]\xi_\alpha{}^j$ dropped), and the cubic term of the last item consists of the commutator terms. The second and third items state that the regulated Feynman-gauge propagator is not invariant under the supercharges: the commutator with $Q^{GF}$ has a component from $A^*$ to $\lambda$, the sum of the compositions through $A$ (the gauge-fixing arrow $A^*\mapsto-2A^*$) and through $\lambda^*$ (the arrow $\lambda^*\mapsto-{\not d}|_{\Sigma^*}\lambda^*$).
\end{remark}

% ANCHOR: end_of_section 7.3
\begin{remark}[Physics notation: special supersymmetry and large distance]
$Q_2\in\Sigma^*$ is a combination of the special supercharges $S^\alpha{}_i$, in $(\mathbf2,\mathbf1,\bar{\mathbf4})$, and $\bar S^{\dot\alpha i}$, in $(\mathbf1,\mathbf2,\mathbf4)$. The multiplication operator $x\mapsto\Gamma(\rho(x)Q_2,-)$ is that of the conformal Killing spinor $\epsilon_\alpha{}^i(x)=-i\,x_m\sigma^m_{\alpha\dot\alpha}\bar\eta^{\dot\alpha i}$, $\bar\epsilon_{\dot\alpha i}(x)=i\,x_m\eta^\alpha{}_i\sigma^m_{\alpha\dot\alpha}$, linear in $x$: the special supersymmetry transformations of the fields are position dependent, and this multiplication operator fails to commute with $H=-\partial^2$. In Remark \ref{rem:c7-ir} the coefficients grow as $|x|$ for $\Delta$ and for $S^\alpha{}_i$, $\bar S^{\dot\alpha i}$, and as $|x|^2$ for $K_b$; the decay $|y|^{-4}$ of a $K_b$ insertion is the borderline of a logarithmic infrared divergence.
\end{remark}

% ANCHOR: after_label prop:c7-select
\begin{remark}[Physics notation: dimensional analysis of the components]
The proposition is dimensional analysis together with ghost-number conservation, gauge invariance and $SO(4)\times SU(4)_R$ covariance. The value of the conformal Killing spinor ($\xi$ at $x_0=0$) has mass dimension $-\tfrac12$ and its Dirac derivative ($4\eta$) dimension $+\tfrac12$, so the transformation under $Q$ has dimension $w_O+\tfrac12$ and the $\eta$-term dimension $w_O-\tfrac12$; a translation parameter has dimension $-1$, rotation and dilatation parameters $0$, and the special conformal parameter $b$ dimension $+1$. In components the conformal Killing spinor \eqref{eq:c8-killing} satisfies
\[
\partial_m\epsilon_\alpha{}^i=-i\,\sigma^m_{\alpha\dot\alpha}\bar\eta^{\dot\alpha i},\qquad \partial_m\bar\epsilon_{\dot\alpha i}=i\,\eta^\alpha{}_i\sigma^m_{\alpha\dot\alpha},\qquad i\,\partial_m\bar\epsilon_{\dot\alpha i}\bar\sigma^{m\dot\alpha\alpha}=4\eta^\alpha{}_i,\qquad -i\,\bar\sigma^{m\dot\alpha\alpha}\partial_m\epsilon_\alpha{}^i=4\bar\eta^{\dot\alpha i},
\]
the last two being the components of ${\not d}(Q_1+\rho(x)Q_2)=4Q_2$.
\end{remark}

% ANCHOR: after_label cor:c7-Q
\begin{remark}[Physics notation: structures of the one-loop supersymmetry variation]
For $O=4\hbar\,w_{1I}w_{2J}\int\omega\operatorname{tr}(\phi_I\phi_J)$ the one-loop transformation under $Q$ is a dimension-$\tfrac52$ combination of structures with the fields $A_I\lambda$ (classical type), $A_m\lambda$ (gauge-variant type) and $c\lambda^*$ (BRST type). In two-component form, with the colour pairing between the two $\mathfrak g$-valued factors, for $Q^i_\alpha$, $Q_1=(\xi_\alpha{}^i,0)$,
\[
(\rho(w)Q_1,\lambda)=w_{ji}\,\xi_\alpha{}^j\varepsilon^{\alpha\beta}\lambda_\beta{}^i,\qquad (\rho(v)Q_1,\lambda)=i\,v_m\bar\lambda_{\dot\alpha i}\bar\sigma^{m\dot\alpha\alpha}\xi_\alpha{}^i,\qquad (Q_1,\lambda^*)=\xi_\alpha{}^i\lambda^{*\alpha}{}_i,
\]
and for $\bar Q_{\dot\alpha i}$, $Q_1=(0,\bar\xi_{\dot\alpha i})$,
\[
(\rho(w)Q_1,\lambda)=-w^{kl}\,\bar\xi_{\dot\alpha k}\varepsilon^{\dot\alpha\dot\beta}\bar\lambda_{\dot\beta l},\qquad (\rho(v)Q_1,\lambda)=i\,v_m\bar\xi_{\dot\alpha i}\bar\sigma^{m\dot\alpha\alpha}\lambda_\alpha{}^i,\qquad (Q_1,\lambda^*)=-\bar\xi_{\dot\alpha i}\bar\lambda^{*\dot\alpha i},
\]
with $w_{ji}\to A_{ji}$, $w^{kl}\to A^{kl}$ for $\rho(A|_{W_{\mathbb R}^\perp})$ and $v_m\to A_m$ for $\rho(A|_V)$; $(\rho(w)Q_1,\lambda)=\iota_w\Gamma(Q_1,\lambda)=\tfrac14w^{kl}\Gamma_{kl}(Q_1,\lambda)$. In $SU(4)_R$ language, $\mathbf4\otimes\bar{\mathbf4}=\mathbf1\oplus\mathbf{15}$, $\wedge^2\mathbf6=\mathbf{15}$ (the adjoint) and $\mathrm{Sym}^2\mathbf6=\mathbf1\oplus\mathbf{20'}$, and $\mathbf{20'}$ does not occur in $\mathbf1\oplus\mathbf{15}$.
\end{remark}

% ANCHOR: after_label cor:c7-S
\begin{remark}[Physics notation: special supercharges at a point]
For $S^\alpha{}_i$ the Killing spinor at $x$ is $\bar\epsilon_{\dot\alpha i}(x)=i\,x_m\eta^\alpha{}_i\sigma^m_{\alpha\dot\alpha}\in S_-\otimes W^*$, and for $\bar S^{\dot\alpha i}$ it is $\epsilon_\alpha{}^i(x)=-i\,x_m\sigma^m_{\alpha\dot\alpha}\bar\eta^{\dot\alpha i}\in S_+\otimes W$. At the position $x$ of the operator the one-loop transformation under $S^\alpha{}_i$ therefore has the structures of the $\bar Q_{\dot\alpha i}$-variation with parameter $i\,x_m\eta^\alpha{}_i\sigma^m_{\alpha\dot\alpha}$, and the one under $\bar S^{\dot\alpha i}$ the structures of the $Q^i_\alpha$-variation with parameter $-i\,x_m\sigma^m_{\alpha\dot\alpha}\bar\eta^{\dot\alpha i}$, with coefficients of their own.
\end{remark}

% ANCHOR: after_label cor:c7-K
\begin{remark}[Physics notation: special conformal transformations at a point]
For a scalar primary of dimension $w$ the classical special conformal transformation is $b^m\big(x^2\partial_m-2x_m\,x\cdot\partial-2w\,x_m\big)=v^m\partial_m+\tfrac w4\operatorname{div}v$. At $x_0$ its data are the translation $v(x_0)$, the local dilatation $\tfrac14\operatorname{div}v(x_0)=-2\,b\cdot x_0$, the local rotation $X_{mn}(x_0)=4(x_{0m}b_n-x_{0n}b_m)$ and the constant second derivatives, linear in $b$. The conformal expectation of Remark \ref{rem:c7-KDelta} is the physics rule that a primary of dimension $w_O+\gamma$ transforms with $-2(w_O+\gamma)\,b\cdot x_0$, the anomalous dimension $\gamma$ of Theorem \ref{thm:c6-dil} entering through the local dilatation.
\end{remark}

% ANCHOR: end_of_section 7.4
\begin{remark}[Physics notation: the generators and their dimensions]
With the dilatation of physics $D=-E=-\tfrac12\Delta$ (the generator $-E$ of [SW, (3)]), the column $\mathrm{ad}_\Delta$ gives
\[
\begin{aligned}&[D,P_b]=P_b,\qquad [D,Q^i_\alpha]=\tfrac12Q^i_\alpha,\qquad [D,\bar Q_{\dot\alpha i}]=\tfrac12\bar Q_{\dot\alpha i},\\ &[D,S^\alpha{}_i]=-\tfrac12S^\alpha{}_i,\qquad [D,\bar S^{\dot\alpha i}]=-\tfrac12\bar S^{\dot\alpha i},\qquad [D,K_b]=-K_b,\end{aligned}
\]
the dimensions $1$, $\tfrac12$, $-\tfrac12$, $-1$ carried by $P$, $Q$, $S$, $K$, while $M$, $\bar M$, $R$ and $D$ have dimension $0$. The column ``weight of output'' is the canonical dimension of the transformed operator, the dimension of $O$ plus that of the generator, and for $S^\alpha{}_i$, $\bar S^{\dot\alpha i}$ evaluated on the Killing spinor $\rho(x)Q_2$ of dimension $-\tfrac12$.
\end{remark}

% ANCHOR: end_of_section 7.5.1
\begin{remark}[Physics notation: anomalies and finite counterterms]
\eqref{eq:c7-defect} is the breaking, by the regulator, of the quantum master equation, i.e. of the Slavnov--Taylor identity of BRST invariance (arity zero) and of the Ward identities of the superconformal generators (arity one); its flow to scale $\infty$ is the one-loop obstruction, a potential anomaly, and Step 2 is the algebraic-renormalisation problem of removing it by finite local counterterms. The finite counterterm of Chapter 4, the Feynman-gauge longitudinal term, reads for $\mathfrak{su}(N)$, with $d^*A|_V=-\partial_mA_m$,
\[
\frac{27}{16}\,\frac{\hbar\operatorname{tr}_{\mathfrak g}(\mathrm{ad}_X\mathrm{ad}_Y)}{16\pi^2(X,Y)}\int d^4x\,(\partial_mA_m,\partial_nA_n)=-\frac{27g^2N}{256\pi^2}\,\hbar\int d^4x\,(\partial_mA^a_m)^2 .
\]
\end{remark}

% ANCHOR: end_of_section 7.5.2
\begin{remark}[Physics notation: Step 3]
Step 3 computes the one-loop corrections to the transformations of the letters and of the bilinears $\operatorname{tr}(\phi_I\phi_J)$ under $Q^i_\alpha$, $\bar Q_{\dot\alpha i}$, $S^\alpha{}_i$, $\bar S^{\dot\alpha i}$: one number per generator for the half-BPS $\mathbf{20'}$, three for the Konishi operator $\operatorname{tr}(\phi_I\phi_I)$, the superconformal primary of the long Konishi multiplet. $S^{(2)}$ is quadratic in the gaugino antifields and contains the coupling of $\Gamma^m(\epsilon_1,\epsilon_2)A_m+\tfrac14\Gamma_{ij}(\epsilon_1,\epsilon_2)A^{ij}$ to $c^*$: in physics language the supersymmetry algebra closes on the fields only up to a field-dependent gauge transformation and the gaugino equation of motion, and the arity-two components express this on-shell closure at one loop.
\end{remark}

% ANCHOR: end_of_section 7.5.3
\begin{remark}[Physics notation: Step 4]
The letters are the scalars $A_I$ (in $SU(4)$ form $A_{ij}$), the field strength $F_A$ with components $F_{mn}$, $D_mA_I$, $[A_I,A_J]$, the gauginos $\lambda_\alpha{}^i$, $\bar\lambda_{\dot\alpha i}$, and their covariant derivatives $D_m=\partial_m+[A_m,-]$. Step 4 is the complete one-loop mixing matrix (the one-loop dilatation operator) on gauge-invariant two-letter operators, in blocks of fixed canonical dimension, ghost number and $SO(4)\times SU(4)_R$ representation. The superficial degree is the power counting of physics: the dimension of the vertex $O$ minus the canonical dimensions of the external lines.
\end{remark}

% ANCHOR: end_of_section 7.5.4
\begin{remark}[Physics notation: Step 5]
Order $\hbar^2$ is two loops, order $g^4$ in the normalisation of Remark \ref{rem:c6-konishi}. The one-loop counterterm vertices inserted into one-loop graphs subtract the subdivergences, $(\log\varepsilon)^2$ is the leading double logarithm of the proper-time cutoff, and the two-loop dilatation read from the $\log\varepsilon$ coefficients is the two-loop anomalous-dimension matrix.
\end{remark}
